\documentclass{resume} % Use the custom resume.cls style

\usepackage[left=0.4 in,top=0.4in,right=0.4 in,bottom=0.4in]{geometry} % Document margins
\usepackage{mathptmx}
\newcommand{\tab}[1]{\hspace{.2667\textwidth}\rlap{#1}} 
\newcommand{\itab}[1]{\hspace{0em}\rlap{#1}}
\name{Khushi Desai} % Your name
% You can merge both of these into a single line, if you do not have a website.
\vspace{-2mm}
\address{+1 (408) 368-0965 \\ New York, NY} 
\address{\href{mailto:khushipdesai@gmail.com}{khushipdesai@gmail.com} \\ \href{https://www.linkedin.com/in/khushidesai1/}{linkedin.com/in/khushidesai1/} \\
\href{https://github.com/khushidesai1}{github.com/khushidesai}}  %

\begin{document}
%----------------------------------------------------------------------------------------
%	EDUCATION
%----------------------------------------------------------------------------------------

\begin{rSection}{Education}
{\bf PhD Computer Science, \textit{Columbia University}} \hfill {\textit{Jan 2025} - \textit{Present}}\\
{\bf MS Computer Science, \textit{Columbia University}} (3.93 GPA) \hfill {\textit{Sept 2023} - \textit{Dec 2024}}\\
\hspace*{6mm} \textbf{Relevant Coursework:} Probabilistic Models and Machine Learning; Deep Learning in Computer Vision; Computer \hspace*{6mm}Graphics; High Performance Machine Learning; Robot Learning \\
{\bf BA Computer Science, \textit{University of California, Berkeley}} (3.64 GPA) \hfill {\textit{Aug 2019} - \textit{May 2023}}\\
\hspace*{6mm} \textbf{Relevant Coursework:} Deep Neural Networks; Machine Learning; Operating Systems; Database Systems
\end{rSection}
\vspace{-2mm}
%----------------------------------------------------------------------------------------
% TECHINICAL STRENGTHS	
%----------------------------------------------------------------------------------------
\begin{rSection}{TECHNICAL SKILLS}
{\bf Languages:} Python, Java, JavaScript, SQL, Bash, Scheme, HTML, CSS, C, RISC-v, GoLang, MongoDB
\\
{\bf Platforms, Frameworks \& APIs:} PyTorch, TensorFlow, SCVI, Scanpy, Anndata, Docker, DynamoDB, Jupyter Notebooks/Lab, React, Redux, AWS, Selenium, Chrome Extensions, Alexa Skill Development
\end{rSection}
\vspace{-2mm}

\begin{rSection}{EXPERIENCE}

\textbf{Graduate Researcher, \textit{Azizi Lab at Columbia University}} \hfill {\it New York, NY} \(|\) {\it Jan 2024 - Present}\\
\textbf{AMICI} 
\textit{Submitted to Nature Biotechnology (\href{https://www.biorxiv.org/content/10.1101/2025.09.22.677860v1}{paper} in review); Presented at MLCB 2024}
\begin{itemize}
    \itemsep -0.5em \vspace{-0.5em}
    \item Developed AMICI, an attention-based neural network architecture to infer cell interactions in single-cell spatial data
    \item Developed semi-synthetic data to validate model performance, and benchmarked against existing SOTA methods 
    \item Interpreted model weights and predictions to analyze results on real-world single-cell Xenium breast cancer data
\end{itemize}
\textbf{Xenium 5k Colorectal Cancer Analysis} \textit{Presented at Single-Cell Genomics 2025}
\begin{itemize}
    \itemsep -0.5em \vspace{-0.5em}
    \item Pre-processing, analyzing and annotating high-throughput single-cell spatial transcriptomics data
    \item Applying AMICI to discover insights into transcriptional targets for colorectal cancer in African American patients
\end{itemize}

\textbf{Undergraduate Researcher, \textit{Berkeley Artificial Intelligence Research Lab}} \hfill {\it Berkeley, CA} \(|\) {\it Sept 2022 - Dec 2023}\\
\textbf{CARFF} \textit{Paper accepted to European Conference on Computer Vision 2024 in Italy (\href{carff.website}{carff.website}, \href{https://link.springer.com/chapter/10.1007/978-3-031-73024-5_14}{paper})}
\begin{itemize}
    \itemsep -0.5em \vspace{-0.5em}
    \item Developed CARFF, an architecture combining VAE, ViTs, and NeRF, to predict probabilistic scenarios and take optimal actions based on occlusions in autonomous vehicle ego-centric scenes
    \item Built a two-stage training process: a conditional autoencoder to encode from partial or complete occlusions in view; and a NeRF decoder to generate scenes from encoded occlusions, enabling safer vehicle decisions
\end{itemize}

\textbf{Software Development Engineer Intern, \textit{Amazon}} \hfill {\it Sunnyvale, CA} \(|\)
{\it May 2023 - Aug 2023}
\begin{itemize}
    \itemsep -0.5em \vspace{-0.5em}
    \item Designed and implemented frontend and backend for a real-time dashboard for Alexa Smart Properties to monitor the set up of devices for hotels, hospitals, etc. 
    % \item Modified on top of existing Amazon React components to build UX of the website
    \item Improved efficiency of Alexa setup for smart property owners through dashboard that will be incorporated on the team's customer facing website
\end{itemize}

\textbf{Software Development Engineer Intern, \textit{Amazon}} \hfill {\it Sunnyvale, CA} \(|\) {\it May 2022 - Aug 2022}
\begin{itemize}
    \itemsep -0.5em \vspace{-0.5em}
    \item Designed and implemented frontend and Java APIs for personalized dashboard for Alexa Voice Services Console, used by Amazon Alexa Certification team to manage third-party client certifications for integrating Alexa into products
    % \item Created React Redux components from scratch to build UX of the website
    % \item Developed Java APIs to pull and transform data from DynamoDB and external databases for certification details
    \item Improved daily productivity for Alexa Certification team by accurately reflecting certification details
\end{itemize}

\textbf{Solutions Architect Intern, \textit{NVIDIA}} \hfill {\it Santa Clara, CA} \(|\) {\it May 2021 - Aug 2021}
\begin{itemize}
    \itemsep -0.5em \vspace{-0.5em}
    \item Developed a Python perception application in a Docker container using PyTorch semantic segmentation, trained on Cityscapes dataset with RTX A6000 GPUs
    \item Built a custom data loader, and trained and evaluated the model on various compute levels, including DGX A100 GPUs
    \item Created a pitch to market NVIDIA SuperPOD using application performance analysis
\end{itemize}

\end{rSection}
\vspace{-2mm}
\begin{rSection}{LEADERSHIP}

\textbf{Instructor, \textit{InspiritAI}} \hfill {\it Cupertino, CA} \(|\) {\it June 2023 - Present}
\begin{itemize}
    \itemsep -0.5em \vspace{-0.5em}
    \item Teaching high school and middle school students to develop AI models using Python libraries and Jupyter Notebooks
\end{itemize}

\textbf{Sponsorship Director, \textit{Cal Hacks}} \hfill {\it Berkeley, CA} \(|\) {\it Feb 2020 - Dec 2021}
\begin{itemize}
    \itemsep -0.5em \vspace{-0.5em}
    \item Raised over \$200,000 for the largest UC Berkeley hackathon by collaborating with 120+ companies, including Microsoft, Google, and Uniswap, with Steve Wozniak as one of our keynote speakers
\end{itemize}
\end{rSection}

\end{document}
